\scanpage{39}
\hfill 40

The results of Chapter I apply. Hence, since $\underline S(X,\underline A)$ is a full subcategory of $\underline P(X,\underline A)$, by I, 3.1, $R_X\circ R_X=R_X$ and by I, 4.12, if $F$ is a sheaf then $R_X(F)=F$. By I, 6.3, $\underline S(X,\underline A)$ is right complete with right roots being $R_X$ applied to the corresponding right root in $\underline P(X,\underline A)$ and if $\underline A$ has a generator, so does $\underline S(X,\underline A)$. Finally, if $\underline A$ is right perfect with a generator, then $\underline S(X,\underline A)$ is abelian, right perfect and has enough injectives.

\figure{\begin{tikzpicture}[scale=.3]\draw (0,0)--(.4,-.03)--(.46,.36);\end{tikzpicture}}
\uline{Proof}:

i) is an immediate consequence of 4.17 and 4.15, Remark ii).

To prove ii) it is sufficient (by 3.1) to construct $R_X(F)$ and $r_F:F\longrightarrow R(\unsure{T|F})$ satisfying the above universal mapping property.

\figure{\begin{tikzpicture}[scale=.35]\draw (0,.03)--(.8,0);\end{tikzpicture}}
Consider $\rho_F:F\longrightarrow P_XF$. By 1.3, $P_XF$ is a sheaf. Let $\{\gamma_i:G_i\longrightarrow P_XF\}$ be the collection of all subsheaves of $P_XF$ through which $\rho_F$ factors; i.e., such that there is a $\delta_i:F\longrightarrow G_i$ with $\gamma_i\unsure{\varepsilon_i|\delta_i}=\rho_F$. Define $R(F)=\bigcap G_i$. Since
\figure{\begin{tikzpicture}[scale=.3]\draw (0,.5)--(.2,0)--(.45,.3);\end{tikzpicture}}
$\underline S(X,\underline A)$ is left closed, and since intersections are left roots (I, § 5, i).) $R(F)$ is a sheaf. Let $\sigma_F:R(F)\longrightarrow P_XF$ be
\mspage{38}
